\section{Introduction}\label{sec:introduction}

A complex manifold is \emph{Brody hyperbolic} if every holomorphic map
$\C\to X$ is constant. On a compact complex manifold, Brody's theorem
identifies this condition with Kobayashi hyperbolicity, the
nondegeneracy of the intrinsic pseudodistance defined by holomorphic
disc chains \cite{Brody1978,Kobayashi1970}. We use ``hyperbolic'' in
this sense throughout.

For a complex manifold of dimension $n$, the canonical bundle
$K_X=\bigwedge^n(T^{1,0}X)^*$ is the line bundle of holomorphic
$n$-forms. A line bundle is \emph{ample} if some positive tensor power
has global sections defining a holomorphic embedding into projective
space. Kobayashi's canonical-ampleness conjecture predicts that the
absence of nonconstant entire curves forces $K_X$ to be ample on a
compact K\"ahler manifold. It therefore connects a restriction on
one-dimensional holomorphic maps with both positivity of top-degree
forms and projective algebraicity. We prove this implication in every
positive dimension.

\begin{theorem}\label{thm:main}
Let $X$ be a compact connected K\"ahler manifold of positive complex
dimension. If every holomorphic map $\C\to X$ is constant, then $K_X$ is
ample. In particular, a positive tensor power of $K_X$ defines a
holomorphic embedding of $X$ into a complex projective space.
\end{theorem}

\subsection{Context and the main difficulty}

Kobayashi's theory supplies the intrinsic metric framework
\cite{Kobayashi1970}; the canonical-ampleness question and its
relationship with curvature are discussed in \cite{Diverio2020}.
Theorem~\ref{thm:main} gives a positive resolution in the smooth
compact K\"ahler category. A major preceding advance is the theorem
of Wu and Yau: negative holomorphic sectional curvature of a K\"ahler metric forces
canonical ampleness on a smooth projective manifold
\cite{WuYau2016}. Tosatti and Yang extended that implication to
compact K\"ahler manifolds without assuming projectivity
\cite{TosattiYang2017}. Diverio--Trapani extended the conclusion to quasi-negative holomorphic
sectional curvature: the curvature is nonpositive everywhere and
strictly negative in every nonzero tangent direction at one point
\cite{DiverioTrapani2019}. Wu--Yau subsequently gave another proof
covering both the negative and quasi-negative cases
\cite{WuYauRemark2016}. Broder--Stanfield allow the negatively curved
metric $\omega$ to be Hermitian and pluriclosed, meaning
$i\partial\bar\partial\omega=0$, while retaining the ambient
K\"ahler hypothesis \cite[Theorem~1.1]{BroderStanfield2026}.
These metric hypotheses give curvature
inequalities unavailable from the absence of entire curves alone.
The present proof begins instead with the derivative bounds supplied
by compact hyperbolicity.

A related line of work uses Gromov's K\"ahler hyperbolicity
\cite{Gromov1991}: a K\"ahler form pulls back to the exterior
derivative of a bounded one-form on the universal covering.
Building on Gromov's general-type result, Chen and Yang proved
canonical ampleness under this hypothesis
\cite[Definitions~2.2 and~2.7, Theorem~2.11]{ChenYang2018}.
This condition is distinct from Brody hyperbolicity
\cite[Corollary~4.4]{ChenYang2018}.

Recent work also establishes a weaker positivity conclusion under the
absence of rational curves. A line bundle is \emph{nef} if, for every
$\eps>0$, it admits a smooth Hermitian metric whose curvature is
bounded below by $-\eps\omega$, where $\omega$ is any fixed background
K\"ahler form. Ou's characterization of uniruled compact K\"ahler manifolds
\cite[Theorem~1.1]{Ou2025} shows that, in the absence of rational
curves, $K_X$ is pseudoeffective: its first Chern class contains a
closed positive current. Ou's theorem also supplies the
lower-dimensional hypothesis in Cao--H\"oring's rational-curve theorem
\cite[Theorem~1.3]{CaoHoring2020}. The latter then shows that failure
of nefness would produce a rational curve. Thus nefness is already
known under this weaker hypothesis. The further issue for ampleness
is positive canonical volume. We retain a direct disc proof of
nefness because it produces bounded canonical potentials and uses
the same Hessian identity as our volume estimate.

Extremal holomorphic discs have long connected intrinsic complex
geometry with analytic estimates, notably in Lempert's theory on
convex domains \cite{Lempert1981}. The extremal problem used here
includes an area penalty on discs in a compact target. Its difficulty
is at the boundary: compact convergence does not by itself justify
variations of an integral over the whole unit disc. We first obtain
boundary Sobolev regularity from maximality. We then construct a
holomorphic frame normalized on the boundary and integrate the needed
infinitesimal sections into actual holomorphic families that keep the
center fixed. These steps make one Hessian calculation available for
two different positivity estimates.

\subsection{The analytic argument}

Put $n=\dim_{\C}X$, let $\D=\{|z|<1\}$, and fix a K\"ahler metric $k$.
A disc $f:\D\to X$ has finite area if the integral of
$k(f'(z),f'(z))$ over $\D$ is finite. Compact hyperbolicity bounds the squared
$k$-norm of the central derivative of every holomorphic disc by a
constant $D$. For a closed
real $(1,1)$-form $q$, evaluated through its associated Hermitian
tensor, use the area and Poisson integrals
\[
 A_q(f)=\frac1\pi\int_\D q(f',f')\,\dd S,
 \qquad
 P_q(f)=\frac2\pi\int_\D\log\frac1{|z|}\,q(f',f')\,\dd S,
\]
where $\dd S$ is Euclidean area. We maximize
\[
 k(f'(0),f'(0))+P_q(f)-\delta A_h(f)
\]
over finite-area discs with a fixed center, where $h$ is a K\"ahler
metric and $\delta>0$. The negative area term gives compactness. Comparing
a maximizing disc with its dilations then gives a linear boundary-area
tail bound. This improves its regularity to continuity on the closed
disc and $W^{1,p}$ for some $p>2$.

The main analytic work is to make boundary-normalized variations at
such an extremum legitimate. A weighted Hardy-space construction, in the tradition of matrix
spectral factorization \cite{WienerMasani1957,HelsonLowdenslager1958}, gives
a holomorphic frame $B=(b_1,\ldots,b_n)$ of the pulled-back tangent bundle with
$B^*hB=I$ on the boundary. We retain enough Sobolev regularity to realize
the sections $zb_j(z)$ as derivatives of actual center-fixed
holomorphic families $F_t$ with $F_0=f$. Let $\cL$ denote the sum of
the complex second derivatives in these $n$ parameter directions,
evaluated at $t=0$. The trace of the
complex Hessian of the area functional is then
\[
 \cL A_h(F_t)=n-A_{R_h}(f),
 \qquad R_h=-\ddc\log\det h,
 \qquad dd^c=i\partial\bar\partial.
\]
Its proof uses an explicit weak boundary limit; it does not assume that
the maximizing disc extends holomorphically across the circle. The
frame and variation constructions apply to boundary-regular
holomorphic discs independently of the global hyperbolicity
hypothesis, which is used to produce those discs.

Two choices of parameters exploit this identity:
\begin{center}
\begin{tabular}{@{}ll@{}}
\toprule
Parameters & Result \\
\midrule
$h=k$, $q=R_k$, $\delta\downarrow0$ & Uniform upper bound for $P_{R_k}$ \\
$q=-h$, $\delta=1$, $R_h\ge-h$ & Uniform lower bound for $h^n/k^n$ \\
\bottomrule
\end{tabular}
\end{center}
The first bound holds on all finite-area discs. The
Poletsky--Rosay disc envelope \cite{Poletsky1991,Rosay2003,DrinovecForstneric2012}
on $\det T^{1,0}X=K_X^*$ with its zero section removed then gives
bounded local plurisubharmonic weights for $K_X$. We include a direct
smoothing argument proving that $K_X$ is nef. Second, $q=-h$ gives the
pointwise estimate
\begin{equation}\label{eq:intro-volume}
 h^n\ge e^{-D}\left(\frac{n}{2n+D}\right)^n k^n
 \quad\text{whenever }R_h\ge-h.
\end{equation}
This estimate is uniform over the indicated metrics, even though the
auxiliary maximizing discs and their regularity constants need not be.

\subsection{From the estimates to ampleness}

Following the Monge--Amp\`ere and volume strategy used by Wu and Yau
\cite{WuYau2016}, the negative-sign Aubin--Yau equation produces K\"ahler metrics $h_t$, for $t>0$,
in the classes $2\pi c_1(K_X)+t[k]$ with $R_{h_t}=-h_t+t k$.
Equation~\eqref{eq:intro-volume} gives a positive lower bound for their
volumes. Passing to the limit in cohomology as $t\downarrow0$ proves
$\int_X c_1(K_X)^n>0$.

The remaining steps use established results with their usual
hypotheses. Holomorphic Morse inequalities make the nef canonical
bundle \emph{big}, meaning that its sections have maximal,
$n$-dimensional asymptotic growth. The resulting sections make $X$
Moishezon, that is, its meromorphic function field has transcendence
degree $n$. A compact K\"ahler Moishezon manifold is projective.
Finally, the absence of rational curves and the log cone theorem
upgrade bigness to ampleness. The final implication is recorded in
\cite[Lemma~2.1]{DiverioTrapani2019}; see also
\cite[Lemma~5.1]{Diverio2020}. All the new estimates are proved
below; the standard existence, envelope, and positivity theorems are
stated at the points where they are used.

\subsection{Organization}

Section~\ref{sec:preliminaries} fixes the curvature, Green-function,
and disc-functional conventions. Section~\ref{sec:extremal-discs}
constructs maximizing discs and proves their boundary regularity.
Section~\ref{var:section} builds the boundary-unitary frames and
centered holomorphic variations, and Section~\ref{sec:hessian}
establishes the determinant and Hessian identities. Their first use,
in Section~\ref{sec:nefness}, produces bounded canonical potentials
and proves nefness; their second use, in Section~\ref{sec:volume},
gives the uniform volume estimate. Section~\ref{sec:positive-intersection}
turns that estimate into positive top self-intersection, and
Section~\ref{sec:ampleness} completes the passage to bigness,
projectivity, and ampleness. Section~\ref{sec:consequences} applies the
result to pluricanonical sections and semialgebraic universal covers.
Appendix~\ref{app:smoothing} gives the bounded-potential smoothing
argument used in Section~\ref{sec:nefness}.
